%%=============================================================================
%% 页眉页脚和脚注
%%=============================================================================

%% 页眉页脚
\RequirePackage{fancyhdr}
\pagestyle{fancy}
\fancyhf{}
\fancypagestyle{xmu@empty}{ %无页眉页脚
	\fancyhf{}}
\fancypagestyle{xmu@headings}{ %有页眉页脚
	\fancyhf{}
	\fancyfoot[C]{\songti\zihao{5} \thepage}

	\fancyhead[CO]{\zihao{5} \songti \nouppercase\leftmark}
	\fancyhead[CE]{\zihao{5} \songti {\xmu@value@title}}
}

%% 脚注样式
% 设置圈型角标
\interfootnotelinepenalty=100000
\RequirePackage{pifont}
\RequirePackage[hang,flushmargin,perpage,symbol*]{footmisc}
\DefineFNsymbols{circled}{{\ding{192}}{\ding{193}}{\ding{194}}
		{\ding{195}}{\ding{196}}{\ding{197}}{\ding{198}}{\ding{199}}{\ding{200}}{\ding{201}}}
\setfnsymbol{circled}

% 设置角标数字的字体和正文一致，并且 1.6em 缩进
\RequirePackage{scrextend}
\deffootnote{1.6em}{0em}{\thefootnotemark \enskip}
